\section{Illustrative Parameterization}
\label{sec:calibration}

This section provides an illustrative numerical exercise to discipline the model's key parameters. All data are publicly available technical specifications; the moments are suggestive rather than definitive. \textbf{These parameters are suggestive and not structurally identified.}

\subsection{Parameter Values}
\label{subsec:parameters}

Table \ref{tab:calibrated_parameters} summarizes the illustrative parameters and their sources.

\begin{table}[htbp]
\centering
\small
\setlength{\tabcolsep}{6pt}
\caption{Illustrative Parameters}
\label{tab:calibrated_parameters}
\begin{tabular}{@{}>{\raggedright\arraybackslash}p{3.2cm}>{\raggedright\arraybackslash}p{2.3cm}>{\raggedright\arraybackslash}p{2.6cm}>{\raggedright\arraybackslash}p{3.5cm}@{}}
\toprule
\textbf{Parameter} & \textbf{Value} & \textbf{Source} & \textbf{Identification} \\
\midrule
$A_\ell$ (attention budget) & 128K tokens & OpenAI API docs & Context window length \\
$K_\ell$ (number of signals) & 50--100 & Multi-agent configs & Typical input dimensions \\
$\gamma$ (attention elasticity) & 0.35 & MMLU + Needle-in-Haystack & Power-law fit $\pm$ 0.05 \\
$\kappa_\ell$ (cost coefficient) & \$0.005 / 1K tokens & OpenAI pricing page & API price per token \\
$\tau^0_{\ell,k}$ (base precision) & 1--10 (heterogeneous) & Signal importance weights & Normalized by $\max_k \tau^0_k$ \\
\bottomrule
\end{tabular}
\end{table}

\subsubsection{Estimation of $\gamma$}

The attention elasticity $\gamma = 0.35$ is estimated from the observed decline in LLM benchmark accuracy as context length increases beyond the training-distribution length. We use three data sources:

\begin{enumerate}
\item \textbf{MMLU long-context variants:} Accuracy scores at context lengths 4K, 8K, 16K, 32K, 64K, 128K tokens. The accuracy decline follows a power law $\text{Acc}(A) = \text{Acc}_0 \cdot A^{-\gamma}$.
\item \textbf{Needle-in-Haystack:} Retrieval accuracy as a function of context length and needle depth. The power-law fit yields $\gamma = 0.32$ (GPT-4) and $\gamma = 0.38$ (Claude-3).
\item \textbf{SWE-bench multi-agent:} Task completion rates for coding pipelines of varying depth, which indirectly identifies the compounding of attention loss.
\end{enumerate}

We estimate the reduced-form elasticity by pooled OLS on the log-linearized regression:
\begin{equation}\label{eq:gamma_regression}
\ln(\text{Acc}_{m,t}) = \alpha_m - \gamma \ln(A_t) + \epsilon_{m,t},
\end{equation}
where $m$ indexes model and $t$ indexes context-length bin. The pooled estimate is $\hat{\gamma} = 0.35$ with a 95\% confidence interval $[0.29, 0.41]$. The between-model standard deviation is 0.06.

\subsubsection{Sensitivity Analysis}

For sensitivity analysis, we report all illustrative predictions under three values: $\gamma = 0.29$ (lower bound), $\gamma = 0.35$ (point estimate), and $\gamma = 0.41$ (upper bound). Table \ref{tab:sensitivity} shows that the optimal depth $L^*$ varies from 3 to 5 layers across this range.

\begin{table}[htbp]
\centering
\caption{Sensitivity of Illustrative Predictions to $\gamma$}
\label{tab:sensitivity}
\begin{tabular}{@{}lccc@{}}
\toprule
\textbf{Outcome} & $\gamma = 0.29$ & $\gamma = 0.35$ & $\gamma = 0.41$ \\
\midrule
Optimal depth $L^*$ & 5 & 4 & 3 \\
Transmission factor $\rho_\ell$ (at $A_\ell = 128$K) & 0.72 & 0.59 & 0.48 \\
Top-layer precision at $L = 5$ (\% of FB) & 68\% & 59\% & 45\% \\
Marginal loss at $L = 5$ (\% per layer) & 4.2\% & 6.8\% & 9.1\% \\
\bottomrule
\end{tabular}
\end{table}

\subsection{Validation}
\label{subsec:validation}

Figure \ref{fig:transmission_factor} plots the endogenous transmission factor $\rho_\ell$ against the attention budget $A_\ell$ for different signal heterogeneity profiles. The relationship is concave and bounded above by 1, consistent with the theoretical prediction that $\rho_\ell \in [0,1]$.

\begin{figure}[htbp]
\centering
\fbox{\parbox{0.8\textwidth}{\centering \vspace{2em} Figure 1: Endogenous Transmission Factor $\rho_\ell$ vs.~Attention Budget $A_\ell$ \vspace{2em}}}
\caption{The transmission factor is concave and bounded above by 1. Higher signal heterogeneity (fewer dominant signals) raises $\rho_\ell$ for any given budget.}
\label{fig:transmission_factor}
\end{figure}

Figure \ref{fig:profit_function} shows the profit function $\Pi(L)$ and the interior optimum $L^*$. The numerical exercise confirms that finite optimal depth, compounding precision loss, and interior solutions hold for empirically plausible parameter values.

\begin{figure}[htbp]
\centering
\fbox{\parbox{0.8\textwidth}{\centering \vspace{2em} Figure 2: Profit Function $\Pi(L)$ and Interior Optimum $L^*$ \vspace{2em}}}
\caption{The profit function is strictly concave in $L$, ensuring a unique finite maximizer. Parameter values: $\gamma = 0.35$, $A_\ell = 128$K, $K_\ell = 50$, $\kappa_\ell = 0.005$.}
\label{fig:profit_function}
\end{figure}

\subsection{Optimal Depth and Attention Budget}
\label{subsec:optimal_depth_budget}

Figure \ref{fig:optimal_depth} plots the optimal chain depth $L^*$ against the attention budget $A_\ell$. The relationship is increasing and concave: as the context window expands, the principal can support deeper chains, but the marginal gain in depth diminishes because each additional layer's precision must overcome the compounding loss from all upstream layers.

\begin{figure}[htbp]
\centering
\fbox{\parbox{0.8\textwidth}{\centering \vspace{2em} Figure 3: Optimal Chain Depth $L^*$ vs.~Attention Budget $A_\ell$ \vspace{2em}}}
\caption{The optimal depth is increasing and concave in the attention budget. Vertical lines mark GPT-3.5 (32K) and GPT-4 (128K) context windows.}
\label{fig:optimal_depth}
\end{figure}

\subsection{Information Loss and Chain Depth}
\label{subsec:info_loss}

Figure \ref{fig:info_loss} plots the cumulative information loss as a function of chain depth. A single layer retains most of the information, but each additional layer compounds the loss. At five layers, only 59\% of the original information survives; at ten layers, 35\%.

\begin{figure}[htbp]
\centering
\fbox{\parbox{0.8\textwidth}{\centering \vspace{2em} Figure 4: Information Retention Rate and Marginal Loss vs.~Chain Depth $L$ \vspace{2em}}}
\caption{A chain of ten AI agents retains only 35\% of the original information. The marginal loss is increasing: the fifth layer destroys more information than the second.}
\label{fig:info_loss}
\end{figure}
